\chapter{L'ABISSO IN FIAMME}\label{l-abisso-di-fiamme}

\hfill\break

Il vento scosse la Grande Casa per tutta la notte, un caldo vento salino
che soffiava dall'Atlantico dopo tre giorni di \emph{sole} fittizio. Lo
sentii anche mentre dormivo: mi fece alzare nel dormiveglia e fu la
colonna sonora di una decina di sogni agitati. Stava ancora bussando
alla finestra dopo il sorgere del Sole, quando mi vestii e andai a
cercare Carol Lawton.

La casa era rimasta senza corrente elettrica per giorni. L'ingresso al
piano di sopra era illuminato debolmente dal bagliore piovoso che
proveniva da una finestra in fondo al corridoio. La scalinata in legno
di quercia scendeva nell'atrio, dove due finestre a belvedere lasciavano
filtrare una luce diurna color rosa pallido. Trovai Carol nel salotto,
stava regolando un antico orologio da tavolo.

«Come sta Diane?» le chiesi.

Carol mi fissò. «Come prima» rispose, dirigendo le sue attenzioni
all'orologio mentre lo caricava con una chiave d'ottone. «Ero da lei un
attimo fa. Non la sto trascurando, Tyler.»

«Non pensavo che lo stessi facendo. E Jason?»

«L'ho aiutato a vestirsi. Sta meglio durante il giorno. Non so perché.
Le notti sono dure, per lui. La notte scorsa è stata\ldots{} dura.»

«Li terrò sott'occhio entrambi.» Non le chiesi neanche se avesse sentito
qualche notiziario, se la FEMA (Federal Emergency Management Agency) o
la Casa Bianca avessero emanato qualche nuova direttiva. Non ce n'era
motivo; l'universo di Carol finiva ai confini della sua proprietà.

«Dovresti riposare un po'.»

«Ho sessantotto anni. Non dormo più come prima. Ma hai ragione, sono
stanca\ldots{} ho bisogno di sdraiarmi. Appena ho finito qui. Questo
orologio resta indietro se non te ne occupi. Tua madre lo regolava tutti
i giorni, lo sapevi questo? E dopo la sua morte Marie lo caricava ogni
volta che faceva le pulizie. Ma Marie ha smesso di venire da circa sei
mesi. Per sei mesi quest'orologio è rimasto bloccato sulle quattro e un
quarto. Come nella vecchia barzelletta, ha segnato l'ora giusta due
volte al giorno.»

«Dovremmo parlare di Jason.» La notte precedente ero troppo esausto per
poter capire bene cosa stesse succedendo: Jason era arrivato senza
preavviso una settimana prima della fine dello Spin e si era ammalato la
notte in cui le stelle erano ricomparse. I suoi sintomi erano
un'intermittente paralisi parziale e la vista annebbiata, oltre alla
febbre. Carol aveva provato a chiedere l'aiuto di un medico ma le
circostanze l'avevano reso impossibile, quindi si era presa cura di lui
da sola, sebbene non fosse stata in grado di diagnosticare il problema
né di somministrargli più di una semplice cura palliativa.

Temeva che stesse per morire. La sua preoccupazione però non si
estendeva al resto del mondo. Jason le aveva detto di non preoccuparsi.
«Presto le cose torneranno alla normalità» aveva detto.

E gli aveva creduto. Il Sole rosso non aveva il potere di terrorizzare
Carol. «Le notti sono dure, comunque» ribadì. Le notti si impossessavano
di Jason come un brutto sogno.

ooo

Controllai per prima Diane.

Carol l'aveva sistemata in una camera al piano superiore - quella che
era stata la sua cameretta, adattata a generica stanza degli ospiti. La
trovai stabile fisicamente, e respirava senza bisogno d'aiuto, ma la
cosa non era del tutto rassicurante. Faceva parte dell'eziologia della
malattia. La marea avanzava e poi si ritirava, ma ogni ciclo portava via
con sé una parte della sua resistenza e della sua forza.

Le poggiai le labbra sulla fronte calda e disidratata, le dissi di stare
a riposo. Non diede alcun segno di avermi sentito.

Poi andai a trovare Jason. Avevo una domanda da fargli.

Da quello che mi aveva detto Carol, Jase era tornato alla Grande Casa
per via di un contrasto alla Perielio. Lei non riusciva a ricordare la
spiegazione di Jason, ma in parte c'entrava suo padre (``E.D. si sta di
nuovo comportando male'' aveva detto Carol) e in parte ``quell'ometto
nero raggrinzito, quello che è morto. Il marziano.''

Il marziano che aveva somministrato a Jason il farmaco della longevità
che lo aveva reso un Quarto. Il farmaco che avrebbe dovuto proteggerlo
da tutto e adesso lo stava uccidendo.

ooo

Era sveglio quando bussai ed entrai nella sua stanza, la stessa che
occupava trent'anni prima, quando eravamo bambini e il mondo girava
attorno a noi e le stelle si trovavano nel posto giusto. C'era un
rettangolo di colore più chiaro nel punto in cui, un tempo, un poster
del sistema solare proteggeva la parete. C'era la moquette, che da
troppo non veniva pulita a vapore e igienizzata chimicamente, dove un
tempo avevamo versato Coca-Cola e sparpagliato briciole in giornate
piovose come questa.

E c'era Jason.

«Questo sembra Tyler» disse.

Giaceva nel letto, vestito - insisteva nel volersi vestire ogni mattina,
aveva detto Carol - indossava pantaloni puliti color cachi e una camicia
di cotone blu. Aveva la schiena appoggiata ai cuscini e sembrava
assolutamente vigile. Borbottai: «Non c'è molta luce qui, Jase.»

«Apri le tendine se vuoi.»

Lo feci, ma entrò solo altra cupa luce color ambra. «Ti dispiace se ti
visito?»

«Certo che no.»

Non guardava me. Stava guardando, sempre che l'inclinazione della testa
significasse qualcosa, una chiazza sulla parete vuota.

«Carol dice che hai avuto problemi di vista.»

«Carol sta attraversando quella che i tuoi colleghi chiamano negazione.
In realtà, sono cieco. Non riesco più a vedere niente da ieri mattina.»

Mi sedetti sul letto accanto a lui. Quando girò la testa verso di me, il
movimento fu fluido ma dolorosamente lento. Presi una torcia a penna
dalla tasca della camicia e la puntai nel suo occhio destro per
controllare se la pupilla si contraeva.

Non lo fece.

Anzi, peggio.

Luccicò. La pupilla luccicò come se vi fossero stati iniettati minuscoli
diamanti.

Jason probabilmente percepì il mio scatto all'indietro.

«Male fino a questo punto?» chiese.

Non fui in grado di rispondere.

Aggiunse, ancora più cupo: «Non posso usare uno specchio. Ti prego, Ty.
Devi dirmi cosa vedi.»

«Questo\ldots{} non so cosa sia, Jason. Non è qualcosa che posso
diagnosticare.»

«Descrivimelo e basta, per favore.»

Provai a mostrare un certo distacco clinico. «È come se negli occhi ti
fossero cresciuti dei cristalli. La sclera pare normale e l'iride non
sembra interessata, ma la pupilla è completamente offuscata da scaglie
di qualcosa di simile alla mica. Non ho mai sentito parlare di niente
del genere. Avrei detto che era impossibile. Non posso curarla.»

Mi allontanai dal letto, trovai una sedia e mi sedetti lì. Per un po'
non ci fu nessun suono a parte il ticchettio dell'orologio sul comodino,
un altro dei pezzi d'antiquariato di Carol in ottime condizioni.

Poi Jason inspirò e si sforzò di fare quello che doveva pensare essere
un sorriso rassicurante. «Grazie. Hai ragione. Non è una patologia che
puoi curare. Ma avrò comunque bisogno del tuo aiuto durante\ldots{} be',
nei prossimi giorni. Carol ci prova, ma è una faccenda molto più grande
di lei.»

«Anche di me.»

Ancora pioggia battente sulla finestra. «L'aiuto di cui ho bisogno non è
soltanto medico.»

«Se hai una spiegazione per tutto questo\ldots»

«Al massimo ne ho una parziale.»

«Allora ti prego condividila con me, Jase, perché comincio ad avere un
po' paura.»

Piegò la testa, come ascoltando un suono che non avevo avvertito o non
potevo sentire. Stette così fino a quando cominciai a chiedermi se si
fosse dimenticato che ero lì. Poi disse: «La versione breve è che il mio
sistema nervoso è stato sopraffatto da qualcosa che va oltre il mio
controllo. Lo stato dei miei occhi è solo una manifestazione esterna.»

«Una malattia?»

«No, ma questo è l'effetto che fa.»

«È contagiosa?»

«Al contrario. Credo sia unica. Una patologia che soltanto io posso
sviluppare - su questo pianeta, almeno.»

«Quindi ha a che fare con la cura per la longevità.»

«In un certo senso sì. Ma io\ldots»

«No, Jase, mi serve una risposta a questo prima che tu dica qualsiasi
altra cosa. La tua condizione - qualunque essa sia - è un risultato
diretto del farmaco che ti ho somministrato?»

«Non un risultato \emph{diretto}, no\ldots{} non è colpa tua in alcun
modo, se è ciò che intendi.»

«Ora come ora non me ne può fregare di meno di chi sia la colpa. Diane è
malata. Non te l'ha detto Carol?»

«Carol mi ha detto qualcosa sull'influenza\ldots»

«Carol ti ha mentito. Non è influenza. È uno stato avanzato di sindrome
da deperimento cardiovascolare. Ho guidato per tremila chilometri
attraverso la fine del mondo o giù di lì, perché sta morendo, Jase, e
c'è solo una cura che mi viene in mente, e tu l'hai appena messa in
dubbio.»

Abbassò di nuovo la testa, forse involontariamente, come se cercasse di
concentrarsi in mezzo a una calca di distrazioni invisibili.

Ma prima che potessi sollecitarlo, disse: «Ci sono degli aspetti della
vita marziana che Wun non ha mai condiviso con te. E.D. l'ha sempre
subodorato e in una certa misura i suoi sospetti erano fondati. Marte ha
prodotto biotecnologie complesse per secoli. Secoli fa, la Quarta Età
era esattamente ciò che Wun ti ha descritto come un trattamento di
longevità e un'istituzione sociale. Ma da allora si è evoluta. Per la
generazione di Wun la Quarta era poco più di una \emph{piattaforma}, un
sistema operativo biologico capace di far funzionare applicazioni
software molto più sofisticate. Non c'è soltanto un quattro, c'è un 4.1,
un 4.2\ldots{} capisci cosa intendo?»

«Quello che ti ho dato\ldots»

«Quello che mi hai dato era il trattamento tradizionale. Un quattro di
base.»

«Ma?»

«Ma\ldots{} da allora l'ho integrato.»

«Anche questa integrazione faceva parte di quello che Wun aveva
trasportato da Marte?»

«Sì. Lo scopo\ldots»

«Lascia perdere lo scopo. Sei assolutamente sicuro di non soffrire degli
effetti del trattamento originale?»

«Sicuro nei limiti del possibile.»

Mi alzai.

Jason mi sentì andare verso la porta. «Posso spiegare,» disse, «e ho
ancora bisogno del tuo aiuto. Prenditi cura di lei con ogni mezzo, Ty.
Spero che sopravviva. Ma ricordati\ldots{} anch'io ho i minuti contati.»

ooo

La cassetta dei medicinali marziani era dove l'avevo nascosta,
indisturbata, dietro il pannello di rivestimento rotto nel seminterrato
della casa di mia madre, e dopo averla presa la portai nella Grande
Casa, attraversando il prato sotto le raffiche di pioggia ambrata.

Carol era in camera di Diane a somministrarle boccate d'ossigeno con la
maschera.

«Dobbiamo usarla con parsimonia,» la rimproverai, «a meno che tu non
possa far comparire un'altra bombola.»

«Le sue labbra erano un po' blu.»

«Fammi vedere.»

Carol si allontanò dalla figlia. Chiusi la valvola e scostai la
maschera. Dovevo fare attenzione all'ossigeno: era indispensabile viste
le sue difficoltà respiratorie, ma allo stesso tempo pericoloso. Se ne
avessi usato troppo, avrebbe potuto romperle gli alveoli polmonari. Ma
se le sue condizioni fossero peggiorate, per mantenere alti i livelli
ematici avrei dovuto somministrarne dosi maggiori, il genere di
ossigenoterapia normalmente fatta con un respiratore automatico. Ma noi
non avevamo un respiratore.

Non avevamo nemmeno i mezzi clinici occorrenti per un'emogas, ma quando
le avevo sfilato la maschera le sue labbra mi erano sembrate abbastanza
nella norma. Il respiro era rapido e superficiale, comunque, e sebbene
avesse aperto gli occhi una volta rimase letargica e inerte.

Quando aprii la cassetta polverosa ed estrassi le fialette marziane e
una siringa ipodermica, Carol mi osservò sospettosa.

«Quello cos'è?»

«Probabilmente l'unica cosa che può salvarle la vita.»

«Davvero? Ne sei sicuro, Tyler?»

Annuii.

«No,» disse, «voglio dire, ne sei \emph{davvero} sicuro? Perché quello è
ciò che hai dato a Jason, no? Quando aveva la Sclerosi Multipla
Atipica.»

Non aveva senso negarlo. «Sì» ammisi.

«Saranno almeno trent'anni che non pratico la medicina, ma non sono una
sprovveduta. Ho fatto una piccola ricerca sulla Sclerosi Multipla
Atipica, dopo l'ultima volta che sei stato qui. Ho sfogliato gli
articoli nelle riviste. E la cosa interessante è che \emph{non esiste}
una cura. Non c'è un farmaco magico. E se ci fosse, non potrebbe essere
trasversalmente specifico per la sindrome da deperimento
cardiovascolare. Quindi ciò che presumo, Tyler, è che stai per
somministrare un agente farmacologico probabilmente connesso a
quell'ometto rugoso che è morto in Florida.»

«Non voglio mettermi a discutere, Carol. È ovvio che hai tratto le tue
conclusioni.»

«Non voglio che tu \emph{discuta}; voglio che mi \emph{rassicuri}.
Voglio che tu mi dica che questo farmaco non farà a Diane ciò che pare
aver fatto a Jason.»

«Non accadrà» ribattei, ma immaginavo che Carol sapesse che stavo
tenendo nascosta di proposito un'avvertenza, il tacito ``per quanto ne
so''.

Studiò il mio volto. «Ci tieni ancora, a lei.»

«Sì.»

«Non smette mai di stupirmi,» sussurrò Carol, «la tenacia dell'amore.»

Infilai l'ago nella vena di Diane.

ooo

A mezzogiorno la casa non era solo calda ma anche così umida che mi
aspettavo che dai soffitti pendesse del muschio. Rimasi seduto vicino a
Diane per assicurarmi che l'iniezione non avesse effetti collaterali
immediati. A un certo punto, bussarono ripetutamente alla porta di casa.
``Ladri,'' pensai, ``saccheggiatori'', ma prima che raggiungessi l'atrio
Carol aveva già aperto e stava ringraziando un uomo corpulento che annuì
e si voltò per andarsene.

«Quello era Emil Hardy» mi spiegò Carol mentre chiudeva la porta. «Ti
ricordi la famiglia Hardy? Erano i proprietari della piccola casa
coloniale sulla Bantam Hill Road. Emil ha stampato un giornale.»

«Un giornale?»

Teneva in mano due fogli di carta spillati. «Emil ha un generatore
elettrico nel garage. Di notte ascolta la radio e prende appunti, poi
stampa un riepilogo e lo consegna nelle case della zona. Questa è la
seconda uscita. È un brav'uomo dalle buone intenzioni. Ma non capisco a
cosa serva leggere certe cose.»

«Posso dargli un'occhiata?»

«Se ti fa piacere.»

Portai i fogli con me, di sopra.

Emil era un giornalista amatoriale encomiabile. Le storie riguardavano
soprattutto crisi avvenute nel Distretto di Columbia e in Virginia - un
elenco ufficiale di zone da evitare e di evacuazioni legate agli
incendi, e tentativi di ripristinare servizi locali. Le sfogliai
rapidamente. Furono un paio di articoli più in basso che catturarono la
mia attenzione.

Il primo era una relazione sulle radiazioni solari misurate di recente a
livello del suolo, che erano aumentate ma neanche lontanamente intense
come previsto. «Gli scienziati del Governo,» spiegava, «sono perplessi
ma abbastanza ottimisti sulle possibilità di una sopravvivenza umana a
lungo termine.» Non veniva citata alcuna fonte, quindi poteva essere
un'invenzione del commentatore o un tentativo di prevenire ulteriore
panico, ma collimava con la mia esperienza fino ad allora: la nuova luce
solare era strana ma non immediatamente fatale.

Nemmeno una parola su come avrebbe impattato sulla resa dei raccolti,
sul clima o sull'ecologia in generale. Né il caldo pestilenziale né
quella pioggia scrosciante sembravano \emph{normali}.

Subito sotto c'era un articolo dal titolo: {LE LUCI DEL CIELO AVVISTATE
	NEL MONDO}.

C'erano le stesse linee a forma di C e di O che mi aveva indicato Simon
in Arizona. Erano state viste a nord fino ad Anchorage e a sud fino a
Mexico City. I servizi che provenivano dall'Europa e dall'Asia erano
frammentari e si occupavano soprattutto della crisi imminente, ma
qualche notizia sulle ``luci'' era giunta anche da lì («Nota,» diceva la
copia di Emil Hardy «i notiziari via cavo disponibili solo a
intermittenza mostrano un video recente dall'India di un fenomeno simile
su larga scala.» Qualunque cosa volesse dire.)

Mentre ero con lei, Diane si svegliò per alcuni attimi.

«Tyler» sussurrò.

Le presi la mano. Era asciutta e stranamente calda.

«Mi dispiace» aggiunse.

«Non hai nulla di cui dispiacerti.»

«Mi dispiace che tu debba vedermi così.»

«Stai migliorando. Ci vorrà del tempo ma starai bene.»

La sua voce era delicata come il suono di una foglia che cade. Si guardò
attorno e riconobbe la camera. I suoi occhi si spalancarono.

«Eccomi qua!»

«Eccoti.»

«Ripeti il mio nome.»

«Diane» dissi. «Diane. Diane.»

ooo

Diane era gravemente malata, ma era Jason che stava morendo. Me lo
confermò quando andai a trovarlo.

Quel giorno non aveva mangiato, me lo aveva detto Carol. Jase aveva
bevuto dell'acqua ghiacciata con la cannuccia, ma aveva rifiutato ogni
altro tipo di liquidi. Riusciva a malapena a muoversi. Quando gli chiesi
di alzare il braccio lo fece, ma con uno sforzo così inimmaginabile e
una velocità così indolente che lo spinsi giù io. Solo la sua voce era
ancora forte, e prevedeva di perdere persino quella: «Se stanotte va
come la notte scorsa non sarò lucido fino all'alba. Domani, chissà?
Voglio parlare finché ne sono in grado.»

«C'è un motivo per cui le tue condizioni peggiorano di notte?»

«Ce n'è uno semplice, credo. Arriveremo anche a quello. Prima voglio che
tu faccia qualcosa per me. La mia valigia era sulla credenza: è ancora
là?»

«È ancora là.»

«Aprila. Ho portato un registratore. Trovamelo.»

Trovai un rettangolo d'acciaio satinato delle dimensioni di un mazzo di
carte da gioco, accanto a una pila di buste manila gialle indirizzate a
nomi che non riconobbi. «È questo?» chiesi, per poi darmi subito
dell'idiota: ovviamente non poteva vedere.

«Se il marchio è Sony, è quello. Sotto dovrebbe esserci una confezione
di memorie nuove.»

«Sì. Trovata.»

«Bene, allora faremo una chiacchierata. Prima che faccia buio, e forse
anche un pochino dopo. E voglio che tu tenga il registratore acceso.
Qualunque cosa accada. Cambia la memoria quando devi farlo, o la
batteria se si sta scaricando. Fallo per me, d'accordo?»

«Finché Diane non ha urgente bisogno di assistenza. Quando vuoi
iniziare?»

Girò la testa. Le pupille incastonate di diamanti dei suoi occhi
brillarono nella strana luce.

«Ora non sarebbe troppo presto.»
